\documentclass[11pt]{article}
\usepackage{amssymb, amsmath}
\usepackage[margin=1in]{geometry}
%\usepackage{geometry}
%\usepackage[dvips]{graphicx}
%\usepackage{draftwatermark}
%\SetWatermarkColor[gray]{0.95}
%\SetWatermarkText{Confidential}
%\SetWatermarkScale{4}
\usepackage{amsmath}
\usepackage{amsthm}
\usepackage{mathtools}
\usepackage{mathrsfs}
\usepackage{amssymb}
\usepackage{verbatim}
\usepackage{algorithmic}
%\usepackage{algorithm2e}
\usepackage[linesnumbered,ruled,vlined]{algorithm2e}
\usepackage{color}
\usepackage[dvipsnames]{xcolor}
%\usepackage{enumerate}
%\usepackage{appendix}
\usepackage{mdframed}
\usepackage{marginnote}
\usepackage{amsfonts}
\usepackage{lineno}

\newcommand{\nota}[1]{ {\tiny \marginnote{#1}[0cm]}}

\newcommand\mydef{\mathrel{\overset{\makebox[0pt]{\mbox{\normalfont\tiny\sffamily def}}}{=}}}
% Keywords command
\providecommand{\keywords}[1]
{
  \small	
  \textbf{\textit{Keywords---}} #1
}

\newcommand{\mon}[1]{\textcolor{magenta}{\textbf{Mon:} #1}}
\newcommand{\Mod}[1]{\ (\mathrm{mod}\ #1)}
\newcommand{\nega}[1]{\bar{#1}}
\newcommand{\GG}{\mathcal{G}}
\newcommand{\qd}{\qed}
\newcommand{\sos}{\text{\sc{SoS}}}
%\newcommand{\sos}{\text{\sc{sos}}}
%\newcommand{\sos}{\text{sos}}
\newcommand{\supp}{suppt}
\newcommand{\y}{y^{N}}
\newcommand{\eps}{\varepsilon}
\newcommand{\n}{\hat{n}}
\newcommand{\uu}{\hat{u}}
\newcommand{\F}{\mathcal{F}}
\newcommand{\M}{\mathcal{M}}
\newcommand{\FF}{\hat{\mathcal{F}}}
\newcommand{\Z}{\mathcal{Z}}
\newcommand{\Cc}{\mathcal{C}}
\newcommand{\D}{\mathcal{D}}
\newcommand{\V}{\mathcal{V}}
\newcommand{\K}{\tilde{K}}
\newcommand{\LAS}{\text{\sc{Las}}}
\newcommand{\las}{\text{\sc{las}}}
\newcommand{\CLP}{\text{\sc{Clp}}}
\newcommand{\clp}{\text{\sc{clp}}}
\newcommand{\SA}{\text{\sc{Sa}}}
\newcommand{\sa}{\text{\sc{sa}}}
\newcommand{\sgn}{\text{sgn }}
\newcommand{\adj}{\text{ adj }}
\newcommand{\rank}{\text{rank }}
\newcommand{\nullity}{\text{nullity }}
\newcommand{\tr}{\text{Tr}}
%\newcommand{\PEXP}[1]{\tilde{\mathbb{E}}{\left [ #1 \right]}}
\newcommand{\PEXP}[1]{L{\left [ #1 \right]}}
\newcommand{\set}[1]{\left\{ #1 \right\}}
\newcommand{\PS}{\mathcal{P}}
\newcommand{\conv}{\textnormal{conv}}
\newcommand{\Ss}{\mathcal{S}}
\newcommand{\X}{\mathcal{X}}
\newcommand{\Aut}{\text{Aut}}
\newcommand{\Ideal}[1]{{\textbf{I}}\left( #1 \right)}
\newcommand{\GIdeal}[1]{\left\langle #1 \right\rangle}
\newcommand{\ca}[1]{\mathcal{#1}}

%%%%%%%%%%%%%%%%% CSP %%%%%%%%%%%%%%%%%%%%%%%%%
\newcommand{\CSP}{\textsc{CSP}}
\newcommand{\IMP}{\textsc{IMP}}
\newcommand{\Pol}{\textsf{Pol}}
\newcommand{\Inv}{\textsf{Inv}}
%\newcommand{\Pol}{\textnormal{Pol}}
\newcommand{\FC}{\textsc{FC}}
\newcommand{\PC}{\textsc{PC}}
\newcommand{\Nsatz}{\textsc{Nsatz}}
\newcommand{\CPC}{\textsc{CPC}}
\newcommand{\CF}{\textsc{CF}}

\newcommand{\Max}{\textsf{Max}}
\newcommand{\Min}{\textsf{Min}}
\newcommand{\Majority}{\textsf{Majority}}
\newcommand{\Minority}{\textsf{Minority }}
\newcommand{\Minorityns}{\textsf{Minority}}
\newcommand{\lex}{\textsf{lex }}
\newcommand{\lexns}{\textsf{lex}}
\newcommand{\grlex}{\textsf{grlex }}
\newcommand{\grlexns}{\textsf{grlex}}
%\newcommand{\grlex}{\textsf{grlex}}
%\newcommand{\mod2}{(\textsf{mod 2})}
\newcommand{\ThB}{\textsc{TH}}
%%%%%%%%%%%%%%%%%%%%%%%%%%%%%%%%%%%%%%%%%%%%%%%%%


\newcommand{\Variety}[1]{{\textbf{V}}\left( #1 \right)}
%\newcommand{\ideal}[1]{\left\langle #1 \right\rangle}
\newcommand{\spn}[1]{\left\langle #1 \right\rangle}
\newcommand{\spns}{\left\langle \Ss \right\rangle}
\newcommand{\Gg}[2]{S_{#1 \setminus #2}}
\newcommand{\cf}[2]{#1_{#2}}
\newcommand{\cone}{\textbf{cone}}
%\newcommand{\sos}{\text{SoS}}
\newcommand{\symsos}{\text{Symmetric-SoS}}
\newcommand{\x}{\mathbf{x}}
\newcommand{\I}{\emph{$\rm{I}$}}
\newcommand{\Zz}{\mathbb{Z}}
\newcommand{\N}{\mathbb{N}}
\newcommand{\1}{\textbf{1}}
\newcommand{\qsos}{$Q\text{-}\sos$ }
\newcommand{\qsost}[1]{Q_{A}(#1)\text{-}\sos }
\newcommand{\qsosAt}[2]{Q_{#1}(#2)\text{-}\sos }
\newcommand{\multideg}{\textnormal{multideg}}
\newcommand{\LM}{\textnormal{LM}}
\newcommand{\LT}{\textnormal{LT}}
\newcommand{\LC}{\textnormal{LC}}
\newcommand{\LCM}{\textnormal{lcm}}
\newcommand{\GB}{\text{Gr\"{o}bner} }
%\newcommand{\reduce}[2]{\overline{#1}^{#2}} - original. Edited below:
\newcommand{\reduce}[2]{{#1}|_{#2}}
\newcommand{\redu}[1]{\overline{#1}}
%\overline{S(f,g)}^G
\newcommand{\paren}[1]{\left ( #1 \right ) }
%margins:
\newcommand{\Field}{\mathbb{Q}}
\newcommand{\Do}{\mathcal{D}}
\newcommand{\Real}{\mathbb{R}}
\newcommand{\Rational}{\mathbb{Q}}
\newcommand{\Complex}{\mathbb{C}}
\newcommand{\Int}{\mathbb{Z}}
\newcommand{\Bool}{\mathbb{B}}
\newcommand{\T}{\mathcal{T}}

\newcommand{\vb}[1]{\mathbf{#1}}

%interpolating polynomials
\newcommand\p{\mathcal{P}}
%interpolating polynomials
\newcommand\q{\mathcal{Q}}

%\newcommand\tt{\mathcal{T}}


\newcommand\bits{{|D|-1}}
\newcommand\uval{{(|D|-1)}}
\newcommand\ideg{???}
\newcommand\db{2(d+k-1)}
\newcommand\ub{2^{2k-1-\ell} \log (2\alpha)}
\newcommand\ds{2^{2k} \log (2\alpha) S}
\newcommand\Cb{2^{2k-1} \log (2\alpha)}
\newcommand\dref{d}
%To repeat main theorems and maintain the same theorem number
% \makeatletter
% \newcommand{\newreptheorem}[2]{\newtheorem*{rep@#1}{\rep@title}\newenvironment{rep#1}[1]{\def\rep@title{#2 \ref*{##1}}\begin{rep@#1}}{\end{rep@#1}}}
% \makeatother

% \newreptheorem{theorem}{Theorem}

% \newtheorem*{theorem*}{Theorem}
% \newtheorem*{corollary*}[theorem*]{Corollary}

\usepackage{lineno}
%\linenumbers
\newcommand*\patchAmsMathEnvironmentForLineno[1]{%
  \expandafter\let\csname old#1\expandafter\endcsname\csname #1\endcsname
  \expandafter\let\csname oldend#1\expandafter\endcsname\csname end#1\endcsname
  \renewenvironment{#1}%
     {\linenomath\csname old#1\endcsname}%
     {\csname oldend#1\endcsname\endlinenomath}}% 
\newcommand*\patchBothAmsMathEnvironmentsForLineno[1]{%
  \patchAmsMathEnvironmentForLineno{#1}%
  \patchAmsMathEnvironmentForLineno{#1*}}%
\AtBeginDocument{%
\patchBothAmsMathEnvironmentsForLineno{equation}%
\patchBothAmsMathEnvironmentsForLineno{align}%
\patchBothAmsMathEnvironmentsForLineno{flalign}%
\patchBothAmsMathEnvironmentsForLineno{alignat}%
\patchBothAmsMathEnvironmentsForLineno{gather}%
\patchBothAmsMathEnvironmentsForLineno{multline}%
}
\usepackage[shortlabels]{enumitem}
\usepackage{url}
\usepackage[colorlinks,
  linkcolor=blue,
  citecolor=blue,%purple,
  filecolor=blue,urlcolor=blue]
{hyperref}
\usepackage[capitalise]{cleveref}
\bibliographystyle{plainurl}% the mandatory bibstyle

\newtheorem{theorem}{Theorem}[section]
\newtheorem{lemma}[theorem]{Lemma}
\newtheorem{definition}[theorem]{Definition}
\newtheorem{proposition}[theorem]{Proposition}
\newtheorem{claim}[theorem]{Claim}
\newtheorem{corollary}[theorem]{Corollary}
\newtheorem{observation}[theorem]{Observation}
\newtheorem{example}[theorem]{Example}
\newtheorem{problem}[theorem]{Problem}
\newtheorem{conjecture}[theorem]{Conjecture}
\newtheorem{remark}[theorem]{Remark}
\newcommand{\ignore}[1]{}
\newcommand{\ale}[1]{\textcolor{red}{\textbf{Alex:} #1}}
\newcommand{\lv}[1]{\textcolor{ForestGreen}{\textbf{Luis:} #1}}
\newcommand{\SOSe}{\rm{SOS}_\varepsilon\text{-complete}}
\usepackage[normalem]{ulem}

\title{On the Degree Automatability of Sum-of-Squares Proofs} %TODO Please add
\date{}
\author{Alex Bortolotti \thanks{University of Applied Sciences and Arts of Southern Switzerland, IDSIA, Lugano, Switzerland. E-mail: \href{mailto:alex.bortolotti@supsi.ch}{\texttt{alex.bortolotti@supsi.ch}}.} \  \and Monaldo Mastrolilli \thanks{University of Applied Sciences and Arts of Southern Switzerland, IDSIA, Lugano, Switzerland. E-mail: \href{mailto:monaldo.mastrolilli@supsi.ch}{\texttt{monaldo.mastrolilli@supsi.ch}}.} \ \and Luis Felipe Vargas\thanks{University of Applied Sciences and Arts of Southern Switzerland, IDSIA, Lugano, Switzerland. E-mail: \href{mailto:luis.vargas@supsi.ch}{\texttt{luis.vargas@supsi.ch}}.}}



%\author{Monaldo Mastrolilli \thanks{{monaldo@idsia.ch}. IDSIA (Istituto Dalle Molle di Studi sull'Intelligenza Artificiale), CH-6962 Lugano-Viganello, Switzerland.}}

\usepackage{fontspec}
